\documentclass[10pt,a4paper]{amsart}
\usepackage[margin=19mm]{geometry}
\usepackage{amsmath,amssymb,mathtools}
\usepackage[scr=rsfs,cal=euler]{mathalfa}
\usepackage{xfrac}
\usepackage[unicode,hidelinks,bookmarks=false]{hyperref}
\newcommand{\ie}{i.e.\ }
\newcommand{\cf}{cf.\ }
\newcommand{\resp}{respectively\ }
\newcommand{\Subsection}{\subsection}
\providecommand{\nobreakdash}{\nobreak\hbox{-}}
\setlength{\emergencystretch}{2em}
\multlinegap0pt
\usepackage{amssymb}
\usepackage{mathtools}
\usepackage{enumerate}
\usepackage{float}
\usepackage{tikz}
\usepackage{dsfont}
\newtheorem{prop}{Proposition}
\newcommand{\ddt}{\frac{d}{dt}}
\newcommand{\norm}[1]{\|#1\|}
\title{A metric boundary theory for Carnot groups}
\author{Nate Fisher}
\date{Source: arXiv:2408.06510v2 (CC BY 4.0)}
\begin{document}
\maketitle

\noindent\emph{Source and licence.} A metric boundary theory for Carnot groups. Version arXiv:2408.06510v2 (CC BY 4.0), distributed by arXiv under the Creative Commons Attribution 4.0 International licence, http://creativecommons.org/licenses/by/4.0/, as recorded in the arXiv licence metadata for this exact version and shown on the abstract page https://arxiv.org/abs/2408.06510v2. Full licence notice: \texttt{licenses/arxiv-2408.06510-CC-BY-4.0.txt}.\par\smallskip

\noindent\textit{Editorial scope.} Counted unit: the article's prop prop:edges (source0.tex lines 928--930). The article's complete original proof of it is delivered with it (source0.tex lines 933--950, one \texttt{proof} environment of 18 lines). No mathematics is written, repaired or re-derived: the statement and the proof are the article's own lines, in the article's own notation, and no retained line is substituted or re-flowed.  Retained uncounted as the source-local context the proof needs: lines 890--895.  Nothing was omitted from any retained span, so the package carries no omission record.  No retained line contains a citation, so the wrapper carries no bibliography and no bibliography entry.  No retained line contains a figure environment, an image-inclusion command or drawing code, so no image is delivered and the package declares no asset dependency.  Editorial formatting only, in addition: an AMS \texttt{amsart} preamble that restates the article's own declarations and macros used by the retained lines, namely the environments \texttt{prop} on separate counters, together with the standard packages those lines need.  The retained environments print in this document as Proposition 1 in order of appearance, whereas the article prints the same items with the numbering of its own full document.  The complete native source of the article as distributed in the arXiv e-print for this exact version is bundled unchanged as \texttt{sources/163696/source0.tex}.
\begin{prop}\label{prop:pansuDiff} Suppose that $p$ is a point on the unit metric sphere $\partial B$ of a polysmooth layered sup norm on $G$ such that $\norm{\pi_i(p)}_{V_i}^{1/i} = 1$ and such that $\norm{\pi_j(p)}_{V_j}^{1/j} < 1$ for $j\neq i$. If $\pi_i(p)$ is in the interior of a facet of the polyhedron $Q_i$ or if $\norm{\cdot}_{V_i}$ is smooth, then the layered sup norm
\[
\norm{x} = \max_{1\leq j \leq s} \left\{\norm{\pi_j(x)}_{V_j}^{1/j}\right \}
\]
is Pansu differentiable at $p$. The principal blow-up or Pansu derivative is a linear function in the coordinates of $x$ coming from the first layer $V_1$ of the stratification.
\end{prop}

\begin{prop}\label{prop:edges} Suppose that $p \in \partial B$, that $\norm{\pi_i(p)}_{V_i}^{1/i} = 1$, and that $\norm{\pi_j(p)}_{V_j}^{1/j} < 1$ for $j\neq i$. Additionally assume that $\norm{\cdot}_{V_i}$ is polyhedral and that $\pi_i(p)$ belongs to a face of $Q_i$ of codimension greater than or equal to 2. In general, the polysmooth layered sup norm
is not Pansu differentiable at $p$, but its principal blow-up is a piecewise-linear function of the coordinates of $x$ coming from the first layer $V_1$ of the stratification.
\end{prop}

\begin{proof}
Following the same arguments as in the proof of Proposition~\ref{prop:pansuDiff}, we have

\begin{align*}
\lim_{t\to 0}\frac{\norm{p\delta_t x} - \norm{p}}{t}
&= \frac{1}{i}\cdot \lim_{t\to0} \;\ddt \norm{\pi_i(p\delta_t x)}_{V_i}.
\end{align*}

Since this norm is polygonal, locally near $\pi_i(p)$, the norm $\norm{\cdot}_{V_i}$ is piecewise-linear where each subdomain is a linear cone.
For a fixed $x \in G$, the curve $\{p\delta_tx\mid t\geq0\}$ is a smooth curve based at $p$ which is polynomial in $t$. Hence for $t$ sufficiently small, the projected curve $\pi_i(p\delta_t x)$ must stay in a single linear cone, say $C_\alpha$ where $\norm{y}_{V_1} = \alpha \cdot y$ for all $y \in C_\alpha$. In fact, we could write down conditions on the coordinates $x$ which would guarantee that $\pi_i(p\delta_tx) \in C_\alpha$ for $t$ sufficiently small. These conditions will depend on the coefficients of the terms which are linear in $t$ in $\pi_i(p\delta_tx)$, and, therefore, will depend only on coordinates of $x$ coming from the first layer of the grading. For such $x \in G$, 

\begin{align*}
\frac{1}{i}\cdot \lim_{t\to0} \;\ddt \norm{\pi_i(p\delta_t x)}_{V_i} &= \frac1i \lim_{t\to 0}\ddt\; \alpha \cdot \pi_i(p\delta_tx)\\
& = \frac1i \;\alpha\cdot \left(\lim_{t\to 0}\;\ddt\; \pi_i(p\delta_tx)\right) = \frac1i \; \alpha \cdot (c_{i_1,1}, \ldots, c_{i_r,1}),
\end{align*}
where again, $c_{i,k}$ is the coefficient of $t^k$ in the $i$th coordinate of $p\delta_t x$. For the same reasons as above is a linear function in the coordinates of the first layer of the stratification. Since this holds for each of the partial functions and linear cones defining the polygonal norm, our principal blow-up function is piecewise-linear and a function of the coordinates of the first layer of the stratification.

\end{proof}

\end{document}
